\documentclass{cut-thesis}
\usepackage[utf8]{inputenc}
\usepackage[T1]{fontenc}
\usepackage{geometry}
\usepackage{graphicx}
\usepackage{float}
\usepackage{hyperref}
\usepackage[english,indonesian]{babel}

\title{Optimasi Pemrosesan Template LaTeX Menggunakan Pendekatan Berbasis AI}
\author{Academic Support Agent}
\date{\today}

\begin{document}
\maketitle

\begin{abstract}

\end{abstract}

\chapter{PENDAHULUAN}


\chapter{TINJAUAN PUSTAKA}


% Original content (commented out):
% % Please refer to frontmatter/titlepage.tex to customize the template
\documentclass{cut-thesis}

% Additional Packages

% Additional commands

\begin{document}

\frontmatter
\input{frontmatter/titlepage}
\input{frontmatter/Acknowledgements}
\input{frontmatter/Abstract}
\input{frontmatter/tocs}


\mainmatter
\input{mainmatter/Chapter-1}
\input{mainmatter/Chapter-2}
\input{mainmatter/Chapter-3}
\input{mainmatter/Chapter-4}
\input{mainmatter/Chapter-5}


\Urlmuskip=0mu plus 1mu\relax\sloppy
\bibliographystyle{IEEEtran}
\bibliography{references}  

\begin{appendices}
\renewcommand{\thechapter}{\Roman{chapter}}
\titleformat{\chapter}[display] {\normalfont\huge\bfseries}{ΠΑΡΑΡΤΗΜΑ~\thechapter}{0pt}{\huge} 

\input{mainmatter/Appendix-1}

\end{appendices}

\printindex 
\end{document}

\end{document}

\end{document}